\subsection{Pre-Registered Ablations}
\label{app:wave1}

Four directed probes on Qwen3-8B-bf16 against \benchL, each with
three stochastic seeds (s1, s2, s3) under a GPT-4o-mini judge. Every
ablation was pre-registered via a kickoff commit with falsifiable
hypotheses before the matching completion commit landed the numbers.
Numbers below are reported under the new $5$-label \DSixID rubric
(\S\ref{sec:eval-method}); pre-registered binary-correctness deltas from
the original commits are retained in parentheses where useful.

\para{Requester-identity strength (Finding 1 refinement).}
Tests whether the social-context-gating result on \DSixDef
collapses or strengthens when the prompt over- or under-specifies the
requester. The ablation runs three identity levels on \DSixID and
\DFiveDef query prompts under an otherwise fixed pipeline:
\textsc{L0} no identity mention, \textsc{L1} the requester's name
only, and \textsc{L2} the requester's name plus social relationship
and conversational context. Under the new $5$-label rubric, $\mathrm{F1}_\mathrm{PU}$
moves $41.9 \to 42.6 \to 34.8$ across \textsc{L0}/\textsc{L1}/\textsc{L2};
ALLOW \textsc{disclose\_correct} rises $47.5\% \to 51.7\% \to 47.5\%$
while DENY leak rises $62.5\% \to 63.8\% \to 72.5\%$. The pre-registered
binary-correctness null ($\Delta = 0.00 \pm 0.87$\,pp on \DSixID,
commit \texttt{9b9c8ba}) survives only at the leak/no-leak granularity:
identity strength tunes disclosure \emph{confidence} (both ALLOW DC
and DENY leak rise together) without producing a real permission
gate. Table~\ref{tab:d6-ablation-sweep} places these three cells
alongside the rest of the sweep.

\para{Stronger-than-BM25 retrieval (Finding 2 retrieval-side stress test).}
Tests whether any retrieval-side upgrade closes the Memobase/MemOS
gap. Three variants are swept on Qwen3-8B: (i) hybrid-RRF fusion of
BM25 and BGE-M3 top-$k$ lists; (ii) an entity-filter that restricts
BM25 to sessions mentioning entities found in the query; (iii) a
temporal-prior that re-weights BGE-M3 scores by recency, with a
$\lambda \in \{0, 0.25, 0.5, 1.0\}$ mini-sweep at s1 followed by
three-seed runs at $\lambda = 0.5$. Hybrid-RRF wins overall:
$\DThreeID$ rises by $+45$\,pp but $\DOneID$ drops by $-22$\,pp;
entity-filter and temporal-prior each trade off differently but no
variant reaches Memobase- or MemOS-level on either \DFourDef or
\DSixDef (commit \texttt{202a7d5}). Under matched $k{=}5$ the
dense-only BGE-M3 baseline at \DThreeID$+37$\,pp, \DOneID$-18$\,pp
is already reported in Table~\ref{tab:p1c-rag}; Table~\ref{tab:wave1-summary}
places the retrieval variants alongside it.

\para{Time-indexed Memobase (scoped null).}
Tests whether giving Memobase a \texttt{query\_time} field on each
retrieval pass helps on cross-session temporal queries. On \benchL
the corpus contains no \texttt{query\_time} metadata, so the filter
degenerates to a no-op; we include the run as a scoped null:
$\DThreeID = 60.6$, $\DSixID = 74.0$, both within the Memobase
baseline range across three seeds (commit \texttt{8036ce3}). This
constrains the claim ``structured memory uses temporal provenance''
to settings where ingest actually stamps time at turn granularity;
we flag this as a future-work prerequisite rather than a null result
on the underlying idea.

\para{Provenance-chunking RAG (Finding 2 \& 3 reinforcement).}
Tests whether prepending a chunk-level header
$\texttt{[session=X day=Y speaker=Z time=T]}$ to each BGE-M3 top-$5$
retrieved passage recovers cross-session reasoning
(Table~\ref{tab:wave1-perdim}). The pre-registered hypotheses were:

\begin{itemize}
\itemsep=1pt\topsep=2pt
\item $\DFourID$ (d8\_temporal) recovers under
  provenance. \textbf{FALSIFIED} ($-24.7$\,pp vs.\ plain BM25 RAG).
\item $\DOneID$ (d5\_cloze) does not recover.
  \textbf{CONFIRMED} ($+4.9$\,pp, flat within seed noise).
\item $\DThreeID$ (d7\_qa) and \DSixID (d4\_permission)
  remain flat. \textbf{FALSIFIED}: $\DThreeID$ collapses by
  $-59.8$\,pp and $\DSixID$ by $-15.0$\,pp.
\end{itemize}

Seven of nine internal dimensions degrade; only d3\_confabulation gains
$+37$\,pp, and the commit body attributes that gain to the underlying
dense retriever surfacing more abstainable context rather than to the
metadata header itself (commit \texttt{6f21ef2}). The combined message
of the retrieval and provenance ablations is that on \benchL no retrieval-side modification ---
neither stronger fusion nor per-chunk provenance tokens --- narrows
the gap left by \DFourID / \DSixID / \DTwoID; the gap is
structural, not a retrieval-strength or metadata-signal gap.

\para{Hypothesised mechanism.}
Two non-mutually-exclusive accounts are consistent with the
direction and per-dim shape of the C1 effect. \emph{(i)~Attention
budget.} Each header expands every retrieved chunk by ${\sim}10$
tokens of high-frequency structured markers ($\texttt{[session=}$,
$\texttt{day=}$, $\texttt{speaker=}$, $\texttt{time=}$); at top-$5$
this is ${\sim}50$ overhead tokens that displace content tokens
under the same context budget. \emph{(ii)~Format-prior shift.}
Structured-syntax markers may activate code-/structured-data
attention regimes rather than narrative-content reading, dampening
multi-hop integration on dimensions that require reasoning over
chunk content (\DFourID, \DSixID, \DThreeID) while sparing
surface-recall on \DOneID. The d3\_confabulation $+37$\,pp gain
under this view is consistent with the format prior also biasing
the reader toward abstention on confabulation-style queries, a
register effect rather than a retrieval-quality effect. We do
not isolate (i) vs.\ (ii) experimentally; the two predict the
same direction for every cell we observe, so the wave-1 evidence
falsifies the metadata-signal hypothesis without identifying the
mechanism by which the header degrades reading.

% Wave-1 ablation summary (H200 §H6 / §H7).
% Populated 2026-04-24 from commits 9b9c8ba (B1), 202a7d5 (B2), 8036ce3 (B3),
% 6f21ef2 (C1). All rows: Qwen3-8B-bf16, GPU 1, 4o-mini judge, 3 seeds.
\begin{table}[h]
\centering
\small
\caption{Pre-registered ablation summary on \benchL. Every row is
Qwen3-8B-bf16 under GPT-4o-mini judge, three stochastic seeds. ``Best
$\Delta$'' and ``worst $\Delta$'' report the largest signed per-dimension
shift (paper-level dim ID in parentheses) relative to plain BM25 RAG at
matched $k{=}5$ on the same reader.}
\label{tab:wave1-summary}
\begin{tabular}{lllll}
\toprule
Ablation & Lever tested & Best $\Delta$ & Worst $\Delta$ & Verdict \\
\midrule
Requester identity strength      & social / prompt  & \multicolumn{2}{c}{$0.00 \pm 0.87$ pp (\DSixID)} & null on \DSixID \\
Hybrid-RRF retrieval             & retriever        & $+45$ pp (\DThreeID) & $-22$ pp (\DOneID) & retrieval-side trade-off \\
Dense BGE-M3 top-5               & retriever        & $+37$ pp (\DThreeID) & $-18$ pp (\DOneID) & similar pattern \\
Time-indexed Memobase            & provenance       & \multicolumn{2}{c}{---}                          & scoped null (no query time) \\
Provenance-chunk header          & provenance       & $+37$ pp (confabulation)\textsuperscript{$\dagger$} & $-70$ pp (counterfactual) & shallow metadata hurts \\
\bottomrule
\end{tabular}

\vspace{2pt}
\footnotesize
\textsuperscript{$\dagger$}Attributable to the dense retriever rather than to the metadata header; see Table~\ref{tab:wave1-perdim}.
\end{table}

% Wave-1 C1 per-dimension delta table.
% Populated 2026-04-24 from commit 6f21ef2.
% Two seed conventions coexist: C1 uses the legacy ablation-runner
% absolute seeds {1, 2, 3} (filenames *_s1_* / *_s2_* / *_s3_*); main-table
% BM25 RAG uses eval.cli convention seeds {1002, 1003, 1004}
% (filenames *_s2_* / *_s3_* / *_s4_*). Both are n=3 stochastic at T=0.3.
\begin{table}[h]
\centering
\small
\caption{Per-dimension accuracy for provenance-chunking RAG on \benchL:
(BGE-M3 top-$5$ with a $\texttt{[session=X day=Y speaker=Z time=T]}$ header
prepended to every retrieved chunk) vs.\ plain BM25 RAG, both with
Qwen3-8B-bf16 reader and GPT-4o-mini judge. Values are three-seed means
$\pm$ standard deviations at $T{=}0.3$, paired against the Qwen3-8B BM25 RAG main-table cell.
Fine-grained source categories are listed against their paper-level
counterparts where applicable; $\Delta$ is reported relative to BM25 RAG.
Seven of nine dimensions degrade; \DOneID is flat within noise;
d3\_confabulation gains $+37$ pp but the commit analysis attributes this
to dense retrieval surfacing more abstainable context, not to the metadata
header itself.}
\label{tab:wave1-perdim}
\begin{tabular}{llrrr}
\toprule
Fine-grained category & Paper dim & BM25 RAG & Provenance RAG & $\Delta$ (pp) \\
\midrule
d1\_conflict       & ---               & $34.3$ & $10.4 \pm 0.3$ & $-23.9$ \\
d2\_anaphora       & ---               & $40.0$ & $14.1 \pm 0.6$ & $-25.9$ \\
d3\_confabulation  & ---               & $24.9$ & $61.6 \pm 0.7$ & $+36.7$ \\
d4\_permission     & \DSixID         & $59.7$ & $44.7 \pm 0.6$ & $-15.0$ \\
d5\_cloze          & \DOneID         & $74.7$ & $79.6 \pm 0.6$ & $+4.9$  \\
d6\_metadata       & \DTwoID         & $48.1$ & $14.4 \pm 0.9$ & $-33.7$ \\
d7\_qa             & \DThreeID       & $84.5$ & $24.7 \pm 0.6$ & $-59.8$ \\
d8\_temporal       & \DFourID        & $34.6$ & $\phantom{0}9.9 \pm 0.0$ & $-24.7$ \\
d10\_counterfactual & ---              & $76.7$ & $\phantom{0}7.1 \pm 1.2$ & $-69.6$ \\
\bottomrule
\end{tabular}
\end{table}


\subsection{Ablation Sweep on \DSixID}
\label{app:retrieval-sweep}

The five-axis ablation sweep referenced in \S\ref{sec:ablation-studies}
covers $142$ \DSixID cells across retrieval-side, reader-scale,
distractor, prompt-side and identity-strength controls. Every cell is
re-judged under the same $5$-label rubric used in
Appendix~\ref{app:dim-permission}. Pooled metrics are reported in
Table~\ref{tab:d6-ablation-sweep}.

\begin{table}[t]
\centering
\caption{Ablation sweep on \DSixID across $142$ cells, pooled across readers and seeds within each row. \textbf{ALLOW-DC}: \%\ of ALLOW items labelled \textsc{disclose\_correct}. \textbf{DENY-leak}: \%\ of DENY items where the leak detector fired. \textbf{Gap}: DENY-leak $-$ ALLOW-DC (positive $=$ anti-policy disclosure). $\mathrm{F1}_\mathrm{PU} = 2 P U / (P+U)$ with $P = 1{-}\mathrm{DENY\,leak}$ and $U = \mathrm{ALLOW\,DC}$.}
\label{tab:d6-ablation-sweep}
\footnotesize
\setlength{\tabcolsep}{4pt}
\begin{tabular}{@{}llrrrrr@{}}
\toprule
\textbf{Lever} & \textbf{Variant} & \textbf{Cells} & \textbf{ALLOW-DC} & \textbf{DENY-leak} & \textbf{Gap} & $\mathrm{F1}_\mathrm{PU}$ \\
\midrule
\multirow{4}{*}{Retrieval} & dense E5 & 15 & 1.2 & 3.0 & +1.8 & 2.4 \\
 & dense BGE-M3 & 14 & 1.1 & 0.9 & -0.2 & 2.1 \\
 & hybrid BM25$+$rerank & 15 & 0.3 & 0.0 & -0.3 & 0.7 \\
 & temporal-prior & 15 & 4.0 & 6.0 & +2.0 & 7.7 \\
\midrule
\multirow{5}{*}{Reader scale} & Qwen3-0.6B & 9 & 31.9 & 37.9 & +6.1 & 42.1 \\
 & Llama-3.2-3B & 9 & 29.2 & 57.1 & +27.9 & 34.7 \\
 & Mistral-7B & 9 & 42.4 & 54.2 & +11.8 & 44.1 \\
 & Qwen3-8B & 9 & 46.6 & 62.6 & +16.1 & 41.5 \\
 & Qwen3-32B & 9 & 53.1 & 79.7 & +26.6 & 29.4 \\
\midrule
\multirow{3}{*}{Distractor} & no distractor (Oracle) & 15 & 40.9 & 57.1 & +16.1 & 41.9 \\
 & $+$ curated distractor & 15 & 39.5 & 59.2 & +19.8 & 40.1 \\
 & $+$ random distractor & 15 & 41.4 & 58.6 & +17.1 & 41.4 \\
\midrule
\multirow{2}{*}{Policy lever} & explicit DENY tag (prompt) & 5 & 30.3 & 40.0 & +9.7 & 40.3 \\
 & provenance-aware reader & 10 & 0.5 & 0.6 & +0.1 & 1.0 \\
\midrule
\multirow{3}{*}{Identity strength} & \textsc{L0} (no identity) & 1 & 47.5 & 62.5 & +15.0 & 41.9 \\
 & \textsc{L1} (name only) & 1 & 51.7 & 63.8 & +12.1 & 42.6 \\
 & \textsc{L2} (name $+$ relation) & 1 & 47.5 & 72.5 & +25.0 & 34.8 \\
\bottomrule
\end{tabular}
\end{table}


The six blocks reinforce Findings~1 and~2 in distinct ways.
\emph{Retrieval} (top block) shows that no upgrade over the BM25 / Oracle
backbone closes the matched-evidence gap on \DSixID: every variant
collapses ALLOW \textsc{disclose\_correct} below $5\%$ and the best
$\mathrm{F1}_\mathrm{PU}$ ($7.7$ for the temporal-prior) is still $5{\times}$
below the lowest Oracle baseline.
\emph{Reader scale} shows the anti-policy gap growing monotonically with
reader capacity ($+6.1$\,pp at Qwen3-0.6B to $+26.6$\,pp at Qwen3-32B);
$\mathrm{F1}_\mathrm{PU}$ peaks in the middle of the reader-scale axis
(Mistral-7B, $44.1$) because the smallest reader trades disclosure for
amnesia and the largest leaks too much.
\emph{Distractor} confirms that Mode~2 anti-policy disclosure is
intrinsic to reader prompt-following rather than a retrieval artefact:
pooled $\mathrm{F1}_\mathrm{PU}$ is $41.9$ / $40.1$ / $41.4$ across the
no-/curated-/random-distractor conditions, and \DFourID accuracy is
flat at $79.1$ / $78.7$ / $79.7\%$. Adding up to $8$K tokens of
irrelevant sessions modulates neither leakage nor cross-session
reasoning.
\emph{Policy lever} contains the blanket-abstention extreme of Mode~1:
the provenance-aware reader has $\mathrm{F1}_\mathrm{PU}{=}1.0$ with $48\%$
\textsc{don't\_know} and $49\%$ \textsc{other} on $200$ items, while the
explicit DENY-tag prompt is roughly neutral on $\mathrm{F1}_\mathrm{PU}$
but shifts $5\%$ of responses to \textsc{refuse} and $26\%$ to
\textsc{don't\_know}.
\emph{Identity strength} (wave-1, Qwen3-8B Oracle) tunes disclosure
\emph{confidence} rather than gating: ALLOW \textsc{disclose\_correct}
rises $47.5\%{\to}51.7\%{\to}47.5\%$ and DENY leak rises
$62.5\%{\to}63.8\%{\to}72.5\%$ across $\textsc{L0}/\textsc{L1}/\textsc{L2}$
identity contexts, with $\mathrm{F1}_\mathrm{PU}$ peaking at $\textsc{L1}$
($42.6$). Stronger requester-identity prompts therefore raise both arms
of the $2{\times}2$ symmetrically, falsifying the hypothesis that
asker-identity cues impose a permission gate at inference time.
Detailed per-reader and per-seed breakdowns are deferred to the project
release; the pooled signal is summarised here.

% Commit trail:
%   4199ba8 B1 kickoff; 9b9c8ba B1 complete (requester-identity null on D6)
%   226abb9 B2 kickoff; 202a7d5 B2 complete (hybrid-RRF wins; trade-offs persist)
%   b1ca7e6 B3 infra;   8036ce3 B3 complete (scoped null, needs query_time)
%   b11d766 C1 kickoff; 6f21ef2 C1 complete (provenance mostly hurts)
% 142-cell sweep judged under 5-label rubric (commit pending) on $2026$-$05$-$04$.
% Output root: MASim/runs/l_20260408_111046/eval_results/ablations/{b1,b2,b3,c1}*/
% Sweep root: out/ablations_l_*/.
